\section{和の中から区分求積の形を取り出す}\label{節:和の区分求積}

連続関数$f$について，区間$[a,b]$を$n$等分し，右端を使って長方形を作ると，その面積の和は
\[
 \frac{b-a}{n}\sum_{k=1}^{n}
 f\left(a+\frac{k(b-a)}{n}\right)
\]
である。$n$を大きくすれば，この和は$\int_a^b f(x)\dd x$へ近づく。とくに$[0,1]$では
\begin{equation}\label{式:単位区間の区分求積}
 \lim_{n\to\infty}\frac1n\sum_{k=1}^{n}f\left(\frac{k}{n}\right)
 =\int_0^1 f(x)\dd x
\end{equation}
となる。

和の極限でこの公式を使うときには，$\frac{k}{n}$が横の位置，$\frac1n$が長方形の幅であることを確認する。分母に$n$が現れるというだけでは足りない。和全体を，幅と高さの積の和へ書き直す必要がある。

\begin{example}{分母をそろえて長方形の幅を見つける}
次の極限を求めよ。
\[
 \lim_{n\to\infty}\sum_{k=1}^{n}\frac{n}{n^2+k^2}.
\]
\end{example}

各項の分母から$n^2$をくくると
\[
 \frac{n}{n^2+k^2}
 =\frac1n\cdot\cfrac{1}{1+\left(\cfrac{k}{n}\right)^2}
\]
となる。したがって，これは$f(x)=\frac1{1+x^2}$の右端を使う区分求積である。この関数は$[0,1]$で連続なので
\[
 \lim_{n\to\infty}\sum_{k=1}^{n}\frac{n}{n^2+k^2}
 =\int_0^1\frac{1}{1+x^2}\dd x.
\]
$x=\tan t$とおけば，$0\leqq t\leqq\frac{\pi}{4}$で$\dd x=(1+\tan^2t)\dd t$だから，右辺は$\frac{\pi}{4}$である。和の一項は$0$へ近づくが，項数も増えるため，和全体は$0$にはならない。

式\eqref{式:単位区間の区分求積}で，右端の代わりに左端や中点を使っても，同じ積分へ収束する。ただし，和に含まれる点の範囲が変わるなら，積分区間も確かめなければならない。例えば
\[
 \frac1n\sum_{k=1}^{2n}f\left(\frac{k}{n}\right)
\]
では，幅は$\frac1n$のままで，点は$0$から$2$まで並ぶ。これは$[0,2]$の区分求積である。項数が$2n$だからといって，幅を勝手に$\frac1{2n}$へ変えてはいけない。

\begin{example}{和の範囲から積分区間を決める}
次の極限を求めよ。
\[
 \lim_{n\to\infty}\sum_{k=n+1}^{2n}\frac1k.
\]
\end{example}

一項を$\frac1k=\frac1n\cfrac1{\cfrac{k}{n}}$と書く。$k=n+1,\ldots,2n$に対応する点は
$1+\frac1n,\ldots,2$であり，$[1,2]$を$n$等分した小区間の右端である。よって
\[
 \lim_{n\to\infty}\sum_{k=n+1}^{2n}\frac1k
 =\int_1^2\frac1x\dd x=\log2.
\]
ここでは$\frac1x$が$x=0$で定義されないことは問題にならない。実際に使った積分区間は$[1,2]$であり，その上で連続だからである。

有限個の端の項を付けたり取り除いたりしても，その差が$0$へ近づくなら極限は変わらない。例えば，連続関数$f$について
\[
 \frac1n\sum_{k=0}^{n}f\left(\frac{k}{n}\right)
 -\frac1n\sum_{k=1}^{n}f\left(\frac{k}{n}\right)
 =\frac{f(0)}n\longrightarrow0
\]
である。しかし，取り除いた項数が$n$とともに増える場合や，端の関数値が大きくなる場合には，この確認を省けない。区分求積の公式を使う範囲を先に決め，残りの差を別に評価するのが確実である。

